\section*{第 \ref{ch:g5:geometry} 章——几何作图}

\begin{solution}{exo:g5:geometry:1}
两种作图均遵循\cref{met:g5:geometry:basic}；
过 $A$ 的垂线与\omterm{def:g4:geometry:def}{平行}线
本身也互相\omterm{def:g4:geometry:def}{垂直}。
\end{solution}

\begin{solution}{exo:g5:geometry:2}
依次作垂线，如\cref{ex:g5:geometry:chain}所示。若
所作的前三个角恰为\omterm{def:g3:shapes:rightangle}{直角}，边长也准确，则
第四个角别无选择，也必为\omterm{def:g3:shapes:rightangle}{直角}——
作得好的长方形会自行检验。
\end{solution}

\begin{solution}{exo:g5:geometry:3}
底边 $[AB]$ 长 $7$ cm；两弧的\omterm{def:g3:shapes:shapes}{半径}分别为 $6$ cm（自 $A$ 画）与 $3$ cm
（自 $B$ 画），交点为 $C$。
\end{solution}

\begin{solution}{exo:g5:geometry:4}
等腰三角形：底边 $5$ cm，两弧均用 $7$ cm 的圆规张口。
等边三角形：底边 $6$ cm，两弧均用 $6$ cm 的圆规张口。
\end{solution}

\begin{solution}{exo:g5:geometry:5}
三角板保证\omterm{def:g3:shapes:rightangle}{直角}；圆规（或直尺）
将 $5$ cm 的边长量至四条边上。
\end{solution}

\begin{solution}{exo:g5:geometry:6}
画边 $[AB]$，长 $4$ cm，任选 $[AD]$ 的方向（不取\omterm{def:g4:geometry:def}{垂直}方向，
以免作出正方形），取 $AD = 4$ cm；
以 $4$ cm 为\omterm{def:g3:shapes:shapes}{半径}，分别从 $B$ 与 $D$ 作弧，交于 $C$。圆规可验证
四条边相等。
\end{solution}

\begin{solution}{exo:g5:geometry:7}
圆心为 $A$、\omterm{def:g3:shapes:shapes}{半径}为 $4$ cm 的圆，以及圆心为
$B$、\omterm{def:g3:shapes:shapes}{半径}为 $3$ cm 的圆，相交于\emph{两}点（一点在
直线 $(AB)$ 上方，另一点在下方），因为 $4 + 3 > 5$，两圆相交。
\end{solution}

\begin{solution}{exo:g5:geometry:8}
一种正确的作法：“1. 画长方形 $ABCD$，其中 $AB = 6$ cm
（底边）及 $BC = 2$ cm。2. 画两段\omterm{def:g3:shapes:shapes}{半径}为 $4$ cm 的弧，
圆心分别为 $D$、$C$（上方两\omterm{def:g4:geometry:polygon}{顶点}），记两弧的交点为
长方形上方的点 $S$。3. 画 $[DS]$ 与 $[CS]$。”（铅笔
尖所画出的图形，正是等腰三角形 $DSC$。）
\end{solution}

\begin{solution}{exo:g5:geometry:9}
以底边两端为圆心，分别以 $3$ cm 与 $4$ cm 为\omterm{def:g3:shapes:shapes}{半径}所作的圆弧，
长 $8$ cm 的底边，永不相交：即便首尾相接，$3 + 4 = 7$ cm 也
无法跨越 $8$ cm。三角形任意两边之和必须始终
大于第三边——否则三角形无法闭合。
\end{solution}

\begin{solution}{exo:g5:geometry:10}
画一条长 $6$ cm 的线段，并取其中点 $O$；过
$O$ 作垂线，在垂线上取两点，距离均为
$2$ cm（分居 $O$ 两侧）；依次连接四个端点。圆规可验证
四条边相等。
\end{solution}

\begin{solution}{exo:g5:geometry:11}
在平面图上，山羊能到达以 $A$ 为圆心的圆盘，
\omterm{def:g3:shapes:shapes}{半径}为 $4$ cm，被距离 $3$ cm 的围栏线截去一块（自 $A$ 量起），其活动区域为
圆盘去掉围栏外侧的弓形部分。沿围栏，她
只能在一条线段上吃草（即圆与围栏线相交之处）；
越过此线段，则无草可吃。
\end{solution}
